% BEGIN REV09_RECTAS27_COORDENADAS
\subsection{Reversible coordinates for ordered pairs of lines}
\label{corpus9:xii:rectas-coordenadas}

Table~\ref{p12-83-c20-s6:tab:h3-27-rectas-u029} fixes the bijection
\(\iota:H_3\to\mathscr L_{27}\), where \(\mathscr L_{27}\) is the
configuration of twenty-seven lines. We retain the central law
\eqref{p12-83-c20-s6:eq:heisenberg-3-nuevo}:
\[
 (a,b,z)(c,d,w)=(a+c,b+d,z+w+bc),\qquad
 (a,b,z)^{-1}=(-a,-b,-z+ab).
\]
Multiplication in both orders verifies the inverse formula: all three
resulting coordinates vanish. All operations in this subsection are
performed in \(\mathbb F_3\).

\begin{proposition}[Reversible encoding of the 729 words]
\label{corpus9:xii:rectas-biyeccion}
For \(w=(\ell_1,h_1,\ell_2,h_2,\ell_3,h_3)\), let
\(L(w)=(\ell_1,\ell_2,\ell_3)\) and
\(H(w)=(h_1,h_2,h_3)\). The map
\begin{equation}
 \mathcal E:\mathbb F_3^6\longrightarrow
                 \mathscr L_{27}\times\mathscr L_{27},\qquad
 \mathcal E(w)=(\iota L(w),\iota H(w))
 \label{corpus9:xii:rectas-codificacion}
\end{equation}
is bijective. In the relative coordinates
\(L=(a,b,z)\), \(R=L^{-1}H=(u,v,t)\), reconstruction is given by
\begin{equation}
 H=(a+u,b+v,z+t+bu),\qquad
 w=(a,a+u,b,b+v,z,z+t+bu).
 \label{corpus9:xii:rectas-reconstruccion}
\end{equation}
\end{proposition}
\begin{proof}
Given an ordered pair \((\ell,h)\) of lines, the table fixes the
unique triples \(L=\iota^{-1}(\ell)\) and
\(H=\iota^{-1}(h)\). Interlacing their coordinates recovers a unique
word \(w\), and evaluation of
\eqref{corpus9:xii:rectas-codificacion} returns the original pair.
Conversely, separating and then interlacing a word preserves all six
coordinates. This establishes both inverse compositions.

The transformation \((L,H)\mapsto(L,L^{-1}H)\) is bijective, with
inverse \((L,R)\mapsto(L,LR)\). The central law gives
\(LR=(a+u,b+v,z+t+bu)\), proving
\eqref{corpus9:xii:rectas-reconstruccion}. In particular, the term
\(bu\) retains the cocycle in the third coordinate. Omitting it changes
the reconstruction by precisely \(bu\); for example,
\(L=(0,1,0)\) and \(R=(1,0,0)\) produce
\(H=(1,1,1)\), whereas componentwise addition produces
\((1,1,0)\).
\end{proof}

The partition in
Theorem~\ref{p12-83-c20-s6:thm:particion-hensel-schlafli-nueva}
is preserved by this bijection: the relative element \(R\) distinguishes
the diagonal, incident, and skew strata, of cardinalities \(27\),
\(270\), and \(432\). Ordered pairs constitute the full domain of
the chart. Adjacency is an additional condition on pairs or their
sequences. Prolongation to complete histories is constructed in
Theorem~\ref{corpus9:xii:rectas-homeomorfismo} through the TPK block
reader.
% END REV09_RECTAS27_COORDENADAS
